\section{The arc Hilbert-Poincar\'{e} series of $y^n=0$ and the Rogers-Ramanujan identities}\label{sec:section_5}

We are now going to compute the Hilbert-Poincar\'{e} series of the focussed arc algebra (at the origin) of the closed subscheme $X$ of $\mathbb{A}^1_k$ given by $y^n=0$, i.e.\ $X$ is the $n$-fold point, where $n\geq 2$. We fix $n$ and as before we denote by $F_i$ and
$f_i$, $i\in \N$, the generators of the defining ideals of $J_{\infty}(X)$ in $k[y_0,y_1,\ldots]$ and of $J_\infty^0(X)$ in $k[y_1,y_2,\ldots]$ respectively. Here we take $F_0:=y_0^n$ and $F_i:=D(F_{i-1})$ for $i\geq 1$, where $D$ is the $k$-derivation that sends $y_i$ to $y_{i+1}$ (see Proposition~\ref{prop:equations_by_deriving}). Then $f_i = F_i|_{y_0=0}$. To describe them explicitly, we need to introduce Bell polynomials. \\

Let $i\geq 1, 1\leq j\leq i$. The \textit{Bell polynomial} $B_{i,j}\in \mathbb{Z}[y_1,\ldots , y_{i-j+1}]$ is defined by the formula
\[  B_{i,j} := \sum \left(\frac{i!}{(1!)^{k_1}(2!)^{k_2}\cdots \bigl((i-j+1)!\bigr)^{k_{i-j+1}}}\right)\, \frac{y_1^{k_1}y_2^{k_2}\cdots y_{i-j+1}^{k_{i-j+1}}}{k_1!k_2!\cdots k_{i-j+1}!},    \]   
where we sum over all tuples $(k_1,k_2, \ldots , k_{i-j+1})$ of nonnegative integers such that 
\[ k_1 + k_2+ \cdots + k_{i-j+1} = j \quad \text{and} \quad k_1 + 2k_2 + \cdots + (i-j+1)k_{i-j+1} = i.  \]
Actually, the coefficient of $y_1^{k_1}\cdots y_{i-j+1}^{k_{i-j+1}}$ equals the number of possibilities to partition a set with $i$ elements into $k_1$ singletons, $k_2$ subsets with two elements, and so on. We put $B_{i,j}:=0$ if $j>i$. 
%We will use the notation
%\[ { i \choose 1_{[k_1]} , 2_{[k_2]},\ldots ,  (i-j+1)_{[k_{i-j+1}]} } := \frac{i!}{(1!)^{k_1}(2!)^{k_2}\cdots \bigl((i-j+1)!\bigr)^{k_{i-j+1}}}  \]  
%for a multinomial coefficient as above. We do not write $k_i$ if $k_i = 1$. 
From the main result of \cite{bruschek_JetAndBell} we deduce:

\begin{Prop}\label{prop:equations_are_bell_polynomials}
We have $F_0 = y_0^n$ and for $i\geq 1$,
\[ F_i  =  \sum_{j=0}^{n-1} \frac{n!}{j!}\, B_{i,n-j}\, y_0^{j}.  \]
It follows that $f_i = 0$ if $i<n$, and for $i\geq n$,
\[f_i  =  n!\, B_{i,n}.\]
\end{Prop}

We endow $k[y_0,y_1,\ldots]$ with the following monomial ordering: for
$\alpha, \beta \in \N^{(\N)}$ we have $y^\alpha > y^\beta$ if and only
if $\wt\, \alpha > \wt\, \beta$ or, in case of equality, the last non-zero
entry of $\alpha - \beta$ is negative (i.e., a weighted reverse lexicographic ordering). The leading term of $F_i$
with respect to this ordering is determined by Proposition~\ref{prop:equations_are_bell_polynomials}:

\begin{Prop}\label{prop:leadingterm_Fi}
Let $i\geq 0$ and write $i = q n + r$ with $0\leq r < n$. The leading term of $F_i$ is 
$$\lt(F_i) = \binom{n}{r} \frac{i!}{(q!)^{n-r}\bigl((q+1)!\bigr)^r}\, y_q^{n-r} y_{q+1}^r. $$
For $i\geq n$, this is also the leading term of $f_i$.
\end{Prop}

It will turn out that these leading terms generate the leading ideal of the ideal $I=(f_i; i\geq n)$ of $k[y_1,y_2,\ldots]$, i.e., the ideal generated by the leading monomials of all polynomials in $I$. Theorem~\ref{thm:hilbert_leading} from the appendix tells us that we can deduce the arc Hilbert-Poincar\'e series from this leading ideal. For this, we need to compute a Gr\"obner basis. All results about Gr\"obner basis theory that we need are collected in the appendix as well.

\begin{Rem}
The results in the appendix are stated for polynomial rings in finitely many variables. In the proof of the next crucial lemma we will use them for countably many variables. We may do this, since we can `approximate' the arc Hilbert-Poincar\'e series according to Proposition~\ref{prop:approximate_by_jet_spaces}. We will explain this more precisely after the proof of the lemma.
\end{Rem}

\begin{Lem}\label{lem:gbbasis_nfold_point}
The leading ideal of $I=(f_i; i\geq n)$ is given by $L(I)=(\lm(f_i); i\geq
n)$. 
\end{Lem}

Before giving the proof of this lemma, we will give some concrete computations for $n=4$ to explain the ideas of the proof. By Gr\"obner basis theory it suffices to show that all $S$-polynomials on the $f_i$ reduce to zero modulo $\{f_i; i\geq n\}$, since the $S_{i,j}$ form a basis of the syzygies on the leading terms of the $f_i$ (see Proposition~\ref{prop:S_polynomials_form_basis} and Theorem~\ref{thm:criterion_for_Groebner_basis}). From Proposition~\ref{prop:leadingterm_Fi} we deduce that
\[  S(f_i,f_j) = S(F_i,F_j)|_{y_0=0}, \]
and so we can equally well show that the $S(F_i,F_j)$ reduce to zero modulo $\{F_i; i\geq 0\}$. Moreover we may restrict to those pairs $F_i,F_j$ for which the leading monomials have a nontrivial common factor by Proposition~\ref{prop:coprime_leading_monomials} (this is Step 2.1 in the proof of the lemma). Let us write the first $F_i$ down for $n=4$:
\begin{align*}
F_0 & = y_0^4 , \\
F_1 & = 4 y_0^3y_1, \\
F_2 & = 12y_0^2y_1^2 + 4 y_0^3 y_2, \\
F_3 & = 24 y_0y_1^3  + 36y_0^2 y_1 y_2 + 4y_0^3 y_3 , \\
F_4 & = 24y_1^4  + 144 y_0y_1^2y_2  + 36y_0^2y_2^2 + 48y_0^2y_1y_3 + 4 y_0^3y_4 , \\
F_5 & = 240 y_1^3 y_2 + 360y_0y_1y_2^2  + 240y_0y_1^2y_3 + 120 y_0^2y_2y_3 + 60 y_0^2y_1y_4 + 4y_0^3y_5 , \\
F_6 & =  1080 y_1^2 y_2^2  + 360 y_0 y_2^3  + 480 y_1^3 y_3  + 1440y_0y_1y_2y_3  + 120 y_0^2y_3^2  + 360 y_0y_1^2y_4 \\
    & \quad  + 180 y_0^2y_2y_4 + 72y_0^2y_1y_5           + 4 y_0^3y_6, \\
F_7 & =   2520y_1 y_2^3   + 5040 y_1^2y_2y_3  + 2520 y_0 y_2^2 y_3+ 1680 y_0y_1y_3^2  +  840 y_1^3y_4 \\
    & \quad + 2520 y_0y_1y_2y_4  + 420 y_0^2y_3y_4 + 504 y_0y_1^2y_5 + 252 y_0^2y_2y_5 + 84 y_0^2y_1y_6 + 4 y_0^3 y_7 .
\end{align*}
We may further reduce the set of $S$-polynomials that have to be checked by invoking Proposition~\ref{prop:reduction_of_syzygy_basis}. For instance, we may forget about $S(F_0,F_3)$ if we have checked that $S(F_0,F_2)$ and $S(F_2,F_3)$ reduce to zero, since $\lm(F_2)$ divides the least common multiple of $\lm(F_0)$ and $\lm(F_3)$. Similarly, using $F_1$, we may forget about $S(F_0,F_2)$. If we do this in a precise way, then we see that we only need to check the following $S$-polynomials between the above $F_i$:
\[ \left\{ \begin{array}{l} S(F_i,F_{i+1}) \text{ for } 0\leq i\leq 6, \\
S(F_1,F_7), S(F_2,F_6) , S(F_3,F_5).
\end{array} \right.  \]
This reduction is explained in Step 1 of the proof.

To see that $S(F_i,F_{i+1})$ reduces to zero, we note the following. We start from 
\begin{equation*}
 \mathcal{R} : 4y_1F_0 - y_0F_1 = 0
\end{equation*}
and we derive this relation. This gives
\begin{equation}\label{eq:S(F_1,F_2)}
 4y_2F_0 + 3y_1F_1 - y_0F_2 = 0,
 \end{equation}
or equivalently,
\[ 12 S(F_1,F_2) = -4y_2F_0. \]
This shows that $S(F_1,F_2)$ reduces to zero modulo $\{F_i; i\in \mathbb{N}\}$. Deriving (\ref{eq:S(F_1,F_2)}) once more gives
 \[ 4y_3F_0 + 7y_2F_1 + 2y_1F_2 - y_0F_3 = 0,\]
or
\[ 24S(F_2,F_3) = - 4y_3F_0 - 7y_2F_1.  \]
Hence, $S(F_2,F_3)$ reduces to zero. Similarly, by deriving the right number of times, we can prove that all $S(F_i,F_{i+1})$ reduce to zero. That is essentially Step 2.2 in the below proof.

Finally, we have to argue why $S(F_1,F_7), S(F_2,F_6)$ and $S(F_3,F_5)$ reduce to zero. That amounts to Step 2.3 in the proof of the lemma. First we derive relation $\mathcal{R}$ four times to find
\[ 4y_5F_0 + 15y_4F_1 + 20y_3F_2 +10y_2F_3 - y_0F_5 = 0.  \]
Note that $F_4$ does not appear here. This equation can be written as
\begin{equation}\label{eq:F_3F_5}
 240S(F_3,F_5) = -4y_5F_0 - 15y_4F_1 - 20y_3F_2
 \end{equation}
and this shows that $S(F_3,F_5)$ reduces to zero. Next we look at $S(F_2,F_6) = \frac{y_2^2}{12}F_2 - \frac{y_0^2}{1080}F_6$. We note that the terms of $1080S(F_2,F_6)$ appear in  
\[ 10y_2 D^3 \mathcal{R} + y_0 D^5\mathcal{R}.  \] 
From this we deduce that
\begin{align*}
 1080S(F_2,F_6) & = - 40 y_2y_4F_0 - 110 y_2 y_3 F_1 - 10 y_1y_2F_3  - 4y_0y_6F_0 - 19y_0y_5 F_1\\
  & \quad  -35 y_0y_4 F_2 -30 y_0y_3 F_3 + y_0y_1F_5. 
\end{align*}
Again $F_4$ does not appear, but this does not yet show that $S(F_2,F_6)$ reduces to zero, since 
\[ \lm(S(F_2,F_6)) = y_0^2y_1^3y_3 < y_0y_1^4y_2 = \lm(y_0y_1F_5)\]
for instance. But we do recognize $240y_1S(F_3,F_5)$ and hence we can replace this using (\ref{eq:F_3F_5}). We find
\begin{align*}
 1080S(F_2,F_6) & = (- 40 y_2y_4 +4y_1y_5 - 4y_0y_6) F_0 + (- 110 y_2 y_3   + 15y_1y_4   - 19y_0y_5) F_1\\
  & \quad  + (20y_1y_3 -35 y_0y_4) F_2 - 30 y_0y_3 F_3. 
\end{align*}
This shows that $S(F_2,F_6)$ reduces to zero. We proceed analogously for $S(F_1,F_7) = \frac{y_2^3}{4}F_1 - \frac{y_0^3}{2520}F_7$. First we look at 
\[ 90 y_2^2 D^2 \mathcal{R} + y_0^2 D^6\mathcal{R}. \]
In there we recognize $2520S(F_1,F_7), 2160y_0y_2S(F_3,F_5)$ and $2160y_1S(F_2,F_6)$. We replace the latter two and we find the following after some computations:
\begin{align*}
 2520S(F_1,F_7) & = (-360y_2^2y_3 + 80 y_1y_2y_4 -8y_1^2y_5-36y_0y_2y_5+ 8y_0y_1y_6- 4y_0^2y_7)F_0 \\
 & \quad + (220 y_1y_2 y_3 - 30y_1^2y_4  -135y_0y_2y_4+ 38y_0y_1y_5 - 23y_0^2y_6) F_1 \\
 & \quad + ( -40y_1^2y_3- 180y_0y_2y_3+70 y_0y_1y_4 - 54y_0^2y_5)F_2 \\
  & \quad  + ( 60 y_0y_1y_3- 65y_0^2y_4 )F_3 - 40y_0^2y_3F_4 .
\end{align*}
Unfortunately this does not show yet that $S(F_1,F_7)$ reduces to zero. We have:
\[ \lm(S(F_1,F_7)) = y_0^3y_1^2y_2y_3 < y_0^2y_1^4y_3 = \lm(y_0^2y_3F_4) = \lm(y_0y_1y_3F_3) = \lm(y_1^2y_3F_2).\]
This implies that $( -40y_1^2y_3 , 60 y_0y_1y_3, - 40y_0^2y_3)$ forms a homogeneous syzygy on the leading terms of $(F_2,F_3,F_4)$. We already know that a basis for these syzygies is given by $S_{2,3}$ and $S_{3,4}$. Indeed, we may compute that
\[ -40y_1^2y_3F_2 + 60 y_0y_1y_3F_3 - 40y_0^2y_3F_4 = -480y_1y_3S(F_2,F_3) + 960y_0y_3S(F_3,F_4)\]
Moreover, we explained that $S(F_2,F_3)$ and $S(F_3,F_4)$ reduce to zero modulo $\{F_i; i\geq 0\}$. Replacing their expressions in terms of the $F_i$, we conclude that $S(F_1,F_7)$ reduces to zero as well.

From this example, we see that it will be useful for the proof of Lemma~\ref{lem:gbbasis_nfold_point} to keep track of the leading monomials of the relevant $S$-polynomials.

\begin{Prop}\label{prop:lm_S_poly_consecutive}
Let $q\geq 1, 0\leq r\leq n-1$. Then
\[ \lm\bigl(S(f_{qn+r} , f_{qn+r+1})\bigr) = \left\{\begin{array}{ll} y_{q-1}y_q^{n-r-2}y_{q+1}^{r+2} & \text{if } q\geq 2, r\neq n-1, \\
y_q^2y_{q+1}^{n-2}y_{q+2} & \text{if } q\geq 2, r = n-1, \\
y_1^{n-r+1}y_2^{r-1}y_3 & \text{if } q=1, r\neq 0.
 \end{array}\right. \]
\end{Prop}

We remark that $S(f_n,f_{n+1})= 0$.

\begin{proof}
In all three cases we have written the second biggest monomial with degree $n+1$ and weight $q(n+1)+r+1$ in the variables $y_1,y_2,\ldots$. We only have to show that the monomial occurs with nonzero coefficient in $S(f_{qn+r} , f_{qn+r+1})$. Using Proposition~\ref{prop:leadingterm_Fi} we see that this $S$-polynomial is a multiple of
\[   (n-r)(q!)(qn+r+1)y_{q+1}f_{qn+r} -(r+1)\bigl((q+1)!\bigr) y_q f_{qn+r+1}. \]
Assume for instance that $q\geq 2$ and $r\neq n-1$. A computation using Proposition~\ref{prop:equations_are_bell_polynomials} gives then 
\[ \frac{(n+r+2)(n!)\bigl( (qn+r+1)!\bigr)}{\bigl((q-1)!\bigr)(q!)^{n-r-3}\bigl((q+1)!\bigr)^{r+1}\bigl((r+2)!\bigr)\bigl((n-r-2)!\bigr)} \neq 0  \]
as the coefficient of $y_{q-1}y_q^{n-r-2}y_{q+1}^{r+2}$ in the above expression. The other cases are treated similarly.
\end{proof}


\begin{Prop}\label{prop:lm_S_poly_difficult}
Let $q\geq 1, 1\leq r\leq n-1$. Then
\[ \lm\bigl(S(f_{qn+r} , f_{(q+1)n+n-r})\bigr) = \left\{\begin{array}{ll} y_{q-1}y_q^{n-r-2}y_{q+1}^{r+1}y_{q+2}^{n-r} & \text{if } q\geq 2, r\neq n-1, \\
y_{q-1}y_{q+1}^{n-2}y_{q+2}^2 & \text{if } q\geq 2, r = n-1, \\
y_1^{n-r}y_2^{r+1}y_3^{n-r-2}y_4 & \text{if } q=1, r\neq n-1, \\
y_1^{2}y_2^{n-2}y_4 & \text{if } q=1, r=n-1.
 \end{array}\right. \]
\end{Prop}

\begin{proof}
Now $S(f_{qn+r} , f_{(q+1)n+n-r})$ is a multiple of
\begin{equation}\label{difficult_S_poly}
 \frac{((q+1)n+n-r)!}{(qn+r)!} y_{q+2}^{n-r}f_{qn+r} - \frac{\bigl((q+2)!\bigr)^{n-r}}{(q!)^{n-r}}y_q^{n-r}f_{(q+1)n+n-r}. 
 \end{equation}
In the first case, we have written the second biggest monomial with degree $2n-r$, weight $(q+1)(2n-r)$, and subject to the additional condition that $y_q$ or $y_{q+2}$ has degree at least $n-r$. It only occurs in the first term of (\ref{difficult_S_poly}) due to the factor $y_{q}^{n-r-2}$. In the second case a small computation using Proposition~\ref{prop:equations_are_bell_polynomials} shows that the second biggest monomial $y_q^2y_{q+1}^{n-3}y_{q+2}^2$ does not occur in $S(f_{qn+r} , f_{(q+1)n+n-r})$ (if $n\geq 3$). We have written the third biggest, which appears with nonzero coefficient in (\ref{difficult_S_poly}). 

If $q=1$, then we can compute that no terms containing only $y_1,y_2,y_3$ occur in (\ref{difficult_S_poly}). We have written the biggest monomial containing $y_4$ of degree $2n-r$, weight $2(2n-r)$, and subject to the additional condition that $y_1$ or $y_3$ has degree at least $n-r$. It appears in (\ref{difficult_S_poly}) with nonzero coefficient.
\end{proof}

\begin{proof}[Proof of Lemma~\ref{lem:gbbasis_nfold_point}]
  We will show that the $f_i$ form a Gr\"obner basis of $I$. Using the notation of the appendix, we will first show that the $S_{i,j}$ with $n\leq i      <j$ and
  \[ \left\{ \begin{array}{lcc} j=i+1 & & \text{\textbf{or}} \\ 
  i = qn + r, j=(q+1)n + n - r \text{ for } q\geq 1, 1\leq r \leq n-1 & & \text{\textbf{or}} \\ 
  f_i\text{ and } f_j \text{ have relatively prime leading monomials } & &  \end{array} \right.  \]
  form a homogeneous basis for the syzygies on the leading terms of the $f_i$. In the second step we will show that all elements of this basis reduce to zero modulo $\{f_i; i\geq n\}$. From Theorem~\ref{thm:criterion_for_Groebner_basis} we conclude then that the $f_i$ are indeed a Gr\"obner basis.\\
  
  \textit{Step 1. The set of $S_{i,j}$ described above forms a basis of the syzygies.} 
  
  From Proposition~\ref{prop:S_polynomials_form_basis} we know already that the set of all $S_{i,j}$ forms a homogeneous basis for the syzygies on the $f_i$. If $\lm(f_i)$ and $\lm(f_j)$ are not coprime then by Proposition~\ref{prop:leadingterm_Fi} the syzygy $S_{i,j}$ is of the type $S_{qn,qn+r}$ for $q\geq 1, 0< r< n$ or $S_{qn+r,qn+s}$, where $q\geq 1, 0< r <n, r<s<2n$. 
  
  For $r$ descending from $n-1$ down to $2$ we use Proposition~\ref{prop:reduction_of_syzygy_basis} with $g_i=f_{qn}, g_j = f_{qn+r}$ and $g_k = f_{qn+r-1}$. Indeed, the least common multiple of $\lm(f_{qn})$ and $\lm(f_{qn+r})$ equals $y_q^ny_{q+1}^r$ and this is of course divisible by $\lm(g_k)=y_q^{n-r+1 }y_{q+1}^{r-1}$. So we remove the syzygies $S_{qn,qn+n-1},\ldots , S_{qn,qn+2}$ and we are still left with a basis.
  
  Next we choose $r\in \{1,\ldots, n-2\}$, and we let $s$ descend from $n$ down to $r+2$. We use again Proposition~\ref{prop:reduction_of_syzygy_basis}, now with $g_i=f_{qn+r}, g_j = f_{qn+s}$ and $g_k = f_{qn+s-1}$. The least common multiple of $\lm(f_{qn+r})$ and $\lm(f_{qn+s})$ equals $y_q^{n-r} y_{q+1}^s$ and this is divisible by $\lm(g_k)=y_q^{n-s+1} y_{q+1}^{s-1}$. For these values of $r$ and $s$ we remove the syzygies $S_{qn+r,qn+s}$ and we still have a basis.
  
  Next we let $r$ go up from 1 to $n-2$ and we choose $s\in \{ n+ 1,n+2,\ldots , 2n-r-1\}$. We take $g_i=f_{qn+r}, g_j = f_{qn+s}$ and $g_k = f_{qn+r+1}$. The least common multiple of $\lm(f_{qn+r})$ and $\lm(f_{qn+s})$ equals then $y_q^{n-r}y_{q+1}^{2n-s}y_{q+2}^{s-n}$. This is divisible by $\lm(g_k)=y_q^{n-r-1}y_{q+1}^{r+1}$. For these values of $r$ and $s$ we can again remove the syzygies $S_{qn+r,qn+s}$ and we keep a basis.
  
  Finally we choose $r\in \{2,\ldots,n-1\}$, and we let $s$ descend from $2n-1$ down to $2n-r+1$. We take $g_i=f_{qn+r}, g_j = f_{qn+s}$ and $g_k = f_{qn+s-1}$. The least common multiple of $\lm(f_{qn+r})$ and $\lm(f_{qn+s})$ equals then $y_q^{n-r}y_{q+1}^{r}y_{q+2}^{s-n}$. This is divisible by $\lm(g_k)=y_{q+1}^{2n-s+1}y_{q+2}^{s-n-1}$. For these values of $r$ and $s$ we remove once more the syzygies $S_{qn+r,qn+s}$ and we find the basis that we were looking for.\\

%\textit{Step 2. Modification of the basis of syzygies.}

%We keep the $S_{i,j}$ with $i$ and $j$ consecutive and the $S_{i,j}$ for which $f_i$ and $f_j$ have relatively prime leading monomial. For $i = qn + %r, j=(q+2)n - r$ with $q\geq 1$ and $r$ descending from $n-1$ to $1$ we replace $S_{i,j}$ by
%\[ T_{i,j}:= \sum_{\alpha = r}^{n-1} \frac{[q!\, (q+2)!]^{\alpha-r}\binom{n}{\alpha} \binom{n-r}{n-\alpha} %(n-\alpha)}{[(q+1)!]^{2\alpha-2r}\binom{n}{r} (n-r) }  y_{q+1}^{\alpha-r}S_{qn+\alpha,(q+2)n-\alpha}.     \]
%Later it will become clear why we choose this expression. Note that the $T_{i,j}$ are homogeneous and that we still have a basis of the syzygies.\\
  
\textit{Step 2. All the elements of this basis reduce to zero modulo $\{f_i; i\geq n\}$.}

\textit{Step 2.1.} First we note that $S(f_i,f_j)$ reduces to zero modulo $\{f_i; i\geq n\}$ if $\lm(f_i)$ and $\lm(f_j)$ are relatively prime by Proposition~\ref{prop:coprime_leading_monomials}.

\vspace{3mm}

\textit{Step 2.2.}  For the other two cases we will exploit the differential structure of the ideal
  $F_0, F_1, \dots$. We have $F_i = D^i(F_0)$, where $D$ is the $k$-derivation determined by $D(y_j) = y_{j+1}$ for $j\geq 0$. Since $F_0 = y_0^n$ and $F_1 = D(y_0^n) = ny_0^{n-1}y_1$ we have the simple
  relation 
  \begin{equation}\label{eq:basicrelation}
    \mathcal{R}\colon ny_1F_0 - y_0F_1=0.
  \end{equation}
Let $q\geq 1$ and $r\in \{0,\ldots, n-1\}$. Applying $D^{q(n+1)+r}$ to the relation $\mathcal{R}$ (using the generalized Leibniz
  rule) and evaluating in $y_0=0$ yields
\begin{align*}
  0 & = - \binom{q(n+1)+r}{n-1} y_{q(n+1)+ r - n + 1} f_{n}\\
  & \quad +  \sum_{\alpha = n}^{q(n+1)+r-1} \binom{q(n+1)+r}{\alpha} \left[n y_{q(n+1)+ r - \alpha + 1} f_{\alpha} -  y_{q(n+1)+ r - \alpha} f_{\alpha+1} \right]\\
  & \quad + n y_{1} f_{q(n+1)+r} \\
    & = n\binom{q(n+1)+r}{q} y_{q+1}f_{qn+r} +
  n\binom{q(n+1)+r}{q-1}y_{q}f_{qn+r+1}\\ & \quad - \left[\binom{q(n+1)+r}{q+1} y_{q+1}f_{qn+r} +
  \binom{q(n+1)+r}{q}y_{q}f_{qn+r+1}\right] + E\\
  & = \frac{(q(n+1)+r)!\, (n-r)}{(q+1)!\, (qn+r)!} y_{q+1}f_{qn+r} - \frac{(q(n+1)+r)!\, (r+1)}{q!\, (qn+r+1)!} y_{q}f_{qn+r+1} + E,
\end{align*}
where we denote by $E$ the remaining terms in the expression of the
derivative. The polynomial $E$ is a $\mathbb{Z}$-linear combination of
$y_{q(n+1)+ r - n + 1} f_{n},\ldots, y_{q+2}f_{qn+r-1}$ and $y_{q-1}f_{qn+r+2},\ldots, y_1f_{q(n+1)+r}$. Note that
$D^{q(n+1)+r}\mathcal{R}$ has weight $q(n+1)+r+1$ and is homogeneous of degree $n+1$
with respect to the standard grading. The monomial $M=y_q^{n-r}y_{q+1}^{r+1}$ is
maximal among those monomials which are of weight $q(n+1)+r+1$ and degree
$n+1$. It cannot appear in $E$ and it is the least common multiple of the leading monomials of $f_{qn+r}$ and $f_{qn+r+1}$. Hence we conclude that 
\[ \frac{(q(n+1)+r)!\, (n-r)}{(q+1)!\, (qn+r)!} y_{q+1}f_{qn+r} - \frac{(q(n+1)+r)!\, (r+1)}{q!\, (qn+r+1)!} y_{q}f_{qn+r+1} \]
is a multiple of the $S$-polynomial $S(f_{qn+r},f_{qn+r+1})$ (this can also easily be deduced from Proposition~\ref{prop:leadingterm_Fi}).

Moreover, we have seen in the proof of Proposition~\ref{prop:lm_S_poly_consecutive} that the second biggest monomial of weight $q(n+1)+r+1$ and degree $n+1$ in $y_1,y_2,\ldots$ does occur in $S(f_{qn+r},f_{qn+r+1})$. Thus the equation
\[  \frac{(q(n+1)+r)!\, (n-r)}{(q+1)!\, (qn+r)!} y_{q+1}f_{qn+r} - \frac{(q(n+1)+r)!\, (r+1)}{q!\, (qn+r+1)!} y_{q}f_{qn+r+1} = - E \]
shows that $S(f_{qn+r},f_{qn+r+1})$ reduces to zero modulo $\{f_i; i\geq n\}$. 

%Let us assume first that $r\neq n-1$. Then the monomial $N = y_{q-1}y_q^{n-r-2}y_{q+1}^{r+2}$ is the second biggest monomial of weight $q(n+1)+r+1$ %and degree $n+1$. We will show that it appears with nonzero coefficient in $E$. Then it has to appear with nonzero coefficient in %$S(F_{qn+r},F_{qn+r+1})$, and the equation $D^{q(n+1)+r}\mathcal{R} = 0$ shows then that $S(F_{qn+r},F_{qn+r+1})$ reduces to zero modulo $\{ F_i ; %i\geq 0\}$. By putting $y_0=0$ we find as well that $S(f_{qn+r},f_{qn+r+1})$ reduces to zero modulo $\{ f_i ; i\geq n\}$ (note that $f_i$ and $F_i$ %have the same leading term by Proposition~\ref{prop:leadingterm_Fi}). In $E$ the monomial $N$ appears only in $y_{q-1}F_{qn+r+2}$, and indeed, by the %generalized Leibniz rule the coefficient of $y_{q-1}F_{qn+r+2}$ in $E$ is
%\[ n\binom{q(n+1)+r}{q-2} - \binom{q(n+1)+r}{q-1} = - \frac{(q(n+1)+r)!\, (n+r+2)}{(q-1)!\, (qn+r+2)!} \neq 0.   \]
%Since we put $y_0=0$ in the end, we are only interested in monomials where $y_0$ does not appear. Thus, for $q=1$ one has to give a similar argument %with $N=y_1^{n-r+1}y_2^{r-1}y_3$, appearing in $y_3F_{n+r-1}$. This fails if $r=0$, but in that case we are dealing with $S(f_n,f_{n+1}) = 0$, so then %there is nothing to prove.

%If $r=n-1$, then the second biggest monomial is $N=y_q^2y_{q+1}^{n-2}y_{q+2}$. This monomial appears in $E$ in $y_{q+2}F_{(q+1)n-2}$ whose coefficient %is
%\[ n\binom{(q+1)(n+1)-2}{q+1} - \binom{(q+1)(n+1)-2}{q+2} =  \frac{((q+1)(n+1)-2)!\, (n+2)}{(q+2)!\, ((q+1)n-2)!} \neq 0.   \]
%\displaybreak[0]

\vspace{3mm}

\textit{Step 2.3.} Now let $q\geq 1, 1\leq r \leq n-1$. We are left with showing that 
\[S(f_{qn+r},f_{(q+1)n + n-r})\] 
reduces to zero modulo $\{f_i ; i\geq n\}$. We use descending induction on $r$, starting with the initial cases $r=n-1$ and $r=n-2$. For $r=n-1$ we consider the relation $\mathcal{R}$ from equation (\ref{eq:basicrelation}). Analogously to Step 2.2, we derive it $(n+1)(q+1) - 1$ times and put $y_0=0$ to find that
\begin{align*}
  0 & = - \binom{(n+1)(q+1) - 1}{n-1} y_{(n+1)(q+1) - n } f_{n}  \\
  & \quad +  \sum_{\alpha = n}^{(n+1)(q+1) - 2} \binom{(n+1)(q+1) - 1}{\alpha} \left[n y_{(n+1)(q+1) - \alpha } f_{\alpha} -  y_{(n+1)(q+1) - 1 - \alpha} f_{\alpha+1} \right]\\
  & \quad + n y_{1} f_{(n+1)(q+1) - 1} \\
    & = \frac{((q+1)(n+1)-1)!\, (n+1)}{(q+2)!\, (qn+n-1)!} y_{q+2}f_{qn+n-1}\\
    & \quad - \frac{((q+1)(n+1)-1)!\, (n+1)}{q!\, (qn+n+1)!} y_{q}f_{qn+n+1} + E,
\end{align*} 
where $E$ is a $\mathbb{Z}$-linear combination of 
\[  y_{(n+1)(q+1) - n } f_{n} , \ldots, y_{q+3}f_{qn+n-2} , y_{q+1} f_{qn+n} , y_{q-1}f_{qn+n+2} , \ldots , y_{1} f_{(n+1)(q+1) - 1}.        \]
However, the coefficient of $y_{q+1} f_{qn+n}$ in $E$ equals
\[   n \binom{(n+1)(q+1) - 1}{qn+n} - \binom{(n+1)(q+1) - 1}{qn+n-1}  = 0.  \]
It follows that
\[ \frac{((q+1)(n+1)-1)!\, (n+1)}{(q+2)!\, (qn+n-1)!} y_{q+2}f_{qn+n-1} - \frac{((q+1)(n+1)-1)!\, (n+1)}{q!\, (qn+n+1)!} y_{q}f_{qn+n+1} \]
is a multiple of $S(f_{qn+n-1},f_{qn + n + 1})$ since the monomial $y_{q}y_{q+1}^{n-1}y_{q+2}$, which is the least common multiple of the leading monomials of $f_{qn+n-1}$ and $f_{qn + n + 1}$, cannot occur in $E$. Moreover, from Proposition~\ref{prop:leadingterm_Fi} and Proposition~\ref{prop:lm_S_poly_difficult} we conclude that $S(f_{qn+n-1},f_{qn + n + 1})$ reduces to zero modulo $\{f_i; i\geq n\}$. 

\vspace{1mm}

Next consider $r=n-2$ (and $n\geq 3$). We look at the two relations
\[ \mathcal{A}_1 : \frac{(q!)^2 (q+2)!}{\bigl((q+1)(n+1) - 2\bigr)!} \, y_{q+2} \, D^{(q+1)(n+1) - 2} \mathcal{R} \, \bigr|_{y_0=0} = 0 \]
and
\[ \mathcal{A}_2 : \frac{ q! \bigl((q+2)!\bigr)^2 }{\bigl((q+1)(n+1)\bigr)!} \,  y_{q} \,  D^{(q+1)(n+1)}\mathcal{R}\, \bigr|_{y_0=0} = 0.    \]
We expand the left hand side of $\mathcal{A}_1$ as a $\mathbb{Q}$-linear combination of
\[ y_{q+2} y_{q(n+1) } f_n , \ldots , y_{q+2} y_{1} f_{q (n+1)+ n-1 },  \]
and the left hand side of $\mathcal{A}_2$ as a $\mathbb{Q}$-linear combination of
\[ y_{q} y_{q(n+1) + 2 } f_n , \ldots , y_{q} y_{1} f_{(q+1)(n+1)  }.  \]
A computation shows that the coefficient of $y_{q+2}^2f_{qn+n-2}$ in $\mathcal{A}_1$ equals
\[   \frac{(n+2) (q!)^2 }{( qn + n - 2)!} \]
and that the coefficient of $y_{q}^2f_{qn+n+2}$ in $\mathcal{A}_2$ is
\[   - \frac{ (n+2) \bigl((q+2)!\bigr)^2 }{( qn + n + 2)!}.    \] 
It follows from Proposition~\ref{prop:leadingterm_Fi} that a multiple of $S(f_{qn+n-2},f_{qn+n+2})$ occurs in the left hand side of $\mathcal{A}_1 + \mathcal{A}_2$. By a similar computation we see that the term of $\mathcal{A}_1$ containing $y_{q+1}y_{q+2}f_{qn+n-1}$ and the term of $\mathcal{A}_2$ containing $y_{q+1}y_{q}f_{qn+n+1}$ form a multiple of $y_{q+1}S(f_{qn+n-1},f_{qn+n+1})$ in $\mathcal{A}_1 + \mathcal{A}_2$. From the above we already know that we may express this as a linear combination of
\begin{center}
$ y_{q+1}y_{(q+1)(n+1) - n } f_{n} , \ldots, y_{q+1}y_{q+3}f_{qn+n-2} ,$\\
$ y_{q+1}y_{q-1}f_{qn+n+2} , \ldots , y_{q+1}y_{1} f_{(q+1)(n+1) - 1}.$   
\end{center}
Putting everything together, we conclude that we can write $S(f_{qn+n-2},f_{qn+n+2})$ as a linear combination of
\begin{center}
 $y_{q+2} y_{q(n+1) } f_n , \ldots , y_{q+2}y_{q+3}f_{qn+n-3},$\\
 $y_{q+2}y_{q}f_{qn+n},  \ldots, y_{q+2} y_{1} f_{q (n+1) + n-1 }, $\\
 $y_{q} y_{q(n+1) + 2 } f_n , \ldots , y_q y_{q+2} f_{qn+n} ,$ \\
 $y_q y_{q-1} f_{qn+n+3}  , \ldots, y_{q} y_{1} f_{(q+1)(n+1)}, $ \\
 $y_{q+1}y_{(q+1)(n+1) - n } f_{n} , \ldots, y_{q+1}y_{q+3}f_{qn+n-2} ,$\\
 $y_{q+1}y_{q-1}f_{qn+n+2} , \ldots , y_{q+1}y_{1} f_{(q+1)(n+1) - 1}.$   
\end{center}
We want to apply Propositions~\ref{prop:lm_S_poly_difficult} and \ref{prop:leadingterm_Fi} to conclude that $S(f_{qn+n-2},f_{qn+n+2})$ reduces to zero modulo $\{f_i;i\geq n\}$. The only problem is the appearance of $y_q y_{q+2} f_{qn+n}$ (twice) in the above list. However, we can compute that its coefficient in $\mathcal{A}_1 + \mathcal{A}_2$ equals zero!

\vspace{1mm}

Finally, let $r\leq n-3$ (and $n\geq 4$). We look at the relations
\[ \mathcal{A}_1 : \frac{(q!)^{n-r} (q+2)!}{\bigl(q(n+1) +r+1 \bigr)!} \, y_{q+2}^{n-r-1} \, D^{q(n+1)+r+1} \mathcal{R} \, \bigr|_{y_0=0} = 0 \]
and
\[ \mathcal{A}_2 : \frac{ q! \bigl((q+2)!\bigr)^{n-r} }{\bigl( q(n+1) + 2n - r - 1 \bigr)!} \,  y_{q}^{n-r-1} \,  D^{q(n+1) + 2n - r - 1 }\mathcal{R}\, \bigr|_{y_0=0} = 0.    \]
We expand the left hand side of $\mathcal{A}_1$ as a $\mathbb{Q}$-linear combination of
\[ y_{q+2}^{n-r-1} y_{ q(n+1) + r + 2 - n } f_n , \ldots , y_{q+2}^{n-r-1} y_{1} f_{q (n+1)+ r + 1},  \]
and the left hand side of $\mathcal{A}_2$ as a $\mathbb{Q}$-linear combination of
\[ y_{q}^{n-r-1} y_{ q(n+1) + n-r } f_n , \ldots , y_{q}^{n-r-1} y_{1} f_{ q(n+1) + 2n - r - 1  }.  \]
As before, we may check that a multiple of 
\[ S(f_{qn+r},f_{(q+1)n+n-r}) \]
occurs in the left hand side of $\mathcal{A}_1 + \mathcal{A}_2$. Similarly, multiples of
\[  y_{q+1}S(f_{qn+r+1},f_{(q+1)n+n-r-1}) \text{ and } y_qy_{q+2}S(f_{qn+r+2},f_{(q+1)n+n-r-2}) \] 
occur there. By induction, we know that the latter two $S$-polynomials reduce to zero modulo $\{f_i;i\geq n\}$ and we replace them by their expression in terms of the $f_i$. More precisely: $S(f_{qn+r},f_{(q+1)n+n-r})$ can be expressed as a linear combination of terms of the form $M f_{a}$ where $M$ is a monomial of degree $n-r$ and where $\wt\, M + a = (q+1)(2n-r)$. The maximum of $\lm(Mf_a)$ can be attained at several places. A careful analysis learns that
\[ \lm(M f_a) \leq y_{q-1}y_q^{n-r-3}y_{q+1}^{r+3}y_{q+2}^{n-r-1} =: N, \] 
where the latter monomial can occur as
\begin{center}
 $\lm(y_{q-1}y_{q+2}^{n-r-1}f_{qn+r+3}), \lm\bigl(y_{q+1}S(f_{qn+r+1},f_{(q+1)n+n-r-1})\bigr)$, \\ 
 or as $\lm\bigl(y_qy_{q+2}S(f_{qn+r+2},f_{(q+1)n+n-r-2})\bigr)$
\end{center}
by Propositions~\ref{prop:leadingterm_Fi} and \ref{prop:lm_S_poly_difficult}. Here (and from now on) we assume that $q\geq 2$. The case $q=1$ can be treated in a similar way. Only the following expressions of the form $Mf_a$ have $N$ as leading monomial:
\begin{center}
 $y_{q+2}^{n-r-1}y_{q-1}f_{qn+r+3}, y_{q-1}y_q^{n-r-3}y_{q+2}^2 f_{(q+1)n+n-r-3}$, \\
 $ y_{q-1}y_q^{n-r-3}y_{q+1}y_{q+2} f_{(q+1)n+n-r-2}, y_{q-1}y_q^{n-r-3}y_{q+1}^2f_{(q+1)n+n-r-1}$.
\end{center}
Since $\lm\bigl(S(f_{qn+r},f_{(q+1)n+n-r})\bigr)< N$ we cannot yet conclude that this $S$-polynomi\-al reduces to zero, but we see that the four expressions above must give rise to a homogeneous syzygy on the leading terms of
\[ f_{qn+r+3},f_{(q+1)n+n-r-3},f_{(q+1)n+n-r-2},f_{(q+1)n+n-r-1}.\]
From Step 1 we know that a basis for these syzygies is given by 
\begin{center}
  $S_{qn+r+3,(q+1)n+n-r-3} , S_{(q+1)n+n-r-3,(q+1)n+n-r-2}$, \\
  and $S_{(q+1)n+n-r-2,(q+1)n+n-r-1}$.
 \end{center}
By induction and by Step 2.2 we know that the corresponding $S$-polynomials reduce to zero modulo $\{f_i;i\geq n\}$. Using this, we conclude that $S(f_{qn+r},f_{(q+1)n+n-r})$ can be expressed as a linear combination of terms $Mf_a$ as above, and with $\lm(Mf_a) < N$. But we may repeat a similar argument to get rid of all monomials between
\[ \lm\bigl(S(f_{qn+r},f_{(q+1)n+n-r})\bigr) = y_{q-1}y_q^{n-r-2}y_{q+1}^{r+1}y_{q+2}^{n-r} \]
and $N$. We just have to remark that at no stage of this process the monomial
\[   y_q^{n-r} y_{q+1}^{r} y_{q+2}^{n-r}  \] 
appears as leading monomial of a term (since in all terms there are factors $y_{q-1}$, $y_{q-2},\ldots$ or $y_{q+3},y_{q+4},\ldots$ involved). This ends the proof of the lemma.
\end{proof}

%Finally we have to treat the syzygies $T_{i,j}$ for $i = qn + r, j=(q+2)n - r$, where $q\geq 1, 1\leq r \leq n-1$. 
%We start from the relation
%\begin{equation}\label{eq:basicrelation2}
%    \mathcal{S}\colon ny_0^{n-r-1}y_1F_0 - y_0^{n-r}F_1=0.
%  \end{equation}
%We will show that the syzygy $T_{i,j}$ reduces to zero by showing that a multiple of
%\[ t_{i,j}:= \sum_{\alpha = r}^{n-1} \frac{[q!\, (q+2)!]^{\alpha-r}\binom{n}{\alpha} \binom{n-r}{n-\alpha} %(n-\alpha)}{[(q+1)!]^{2\alpha-2r}\binom{n}{r} (n-r) }  y_{q+1}^{\alpha-r}S(F_{qn+\alpha},F_{(q+2)n-\alpha}) \]
%appears in $D^{(q+1)(2n-r)-1}\mathcal{S}$. For this, we have to investigate the coefficients of $y_{q+1}^{\alpha - r}y_{q+2}^{n-\alpha}F_{qn+\alpha}$ %and of $y_{q+1}^{\alpha - r}y_{q}^{n-\alpha}F_{(q+2)n-\alpha}$ in $D^{(q+1)(2n-r)-1}\mathcal{S}$, for $r\leq \alpha \leq n-1$. Using the generalized %Leibniz rule once more, we see that the first one equals
%\begin{align*}
%& n\binom{n-r-1}{n-\alpha}\binom{(q+1)(2n-r)-1}{q , (q+1)_{(\alpha-r-1)} , (q+2)_{(n-\alpha)} , qn + \alpha  } \\
%& \quad + n\binom{n-r-1}{\alpha - r}\binom{(q+1)(2n-r)-1}{(q+1)_{(\alpha-r+1)} , (q+2)_{(n-\alpha - 1 )} , qn + \alpha  } \\
%& \quad - \binom{n-r}{n-\alpha}\binom{(q+1)(2n-r)-1}{ (q+1)_{(\alpha-r)} , (q+2)_{(n-\alpha)} , qn + \alpha -1 } .
%\end{align*}
%This can be simplified to
%\[ \frac{((q+1)(2n-r)-1)!\,\binom{n-r}{n-\alpha}\, (2n-r)\,(n - \alpha) }{ \bigl((q+1)!\bigr)^{\alpha - r} \, \bigl((q+2)!\bigr)^{n - \alpha} \,(q n + %\alpha)! \,(n-r)      }    . \]
%Similarly, we find
%\[ - \frac{((q+1)(2n-r)-1)!\,\binom{n-r}{n-\alpha}\, (2n-r)\,(n - \alpha) }{ \bigl((q+1)!\bigr)^{\alpha - r}\,  (q!)^{n - \alpha}\, ((q+2) n- %\alpha)!\, (n-r)         }   \] 
%for the coefficient of $y_{q+1}^{\alpha - r}y_{q}^{n-\alpha}F_{(q+2)n-\alpha}$. It follows from Proposition~\ref{prop:leadingterm_Fi} that
%\[ \frac{((q+1)(2n-r)-1)!\,\binom{n-r}{n-\alpha}\,\binom{n}{\alpha}\, (2n-r)\,(n - \alpha) }{(q!)^{n-\alpha}\, \bigl((q+1)!\bigr)^{2\alpha - r} \, %\bigl((q+2)!\bigr)^{n - \alpha}  \,(n-r)} y_{q+1}^{\alpha-r}S(F_{qn+\alpha},F_{(q+2)n-\alpha}) \]
%appears in $D^{(q+1)(2n-r)-1}\mathcal{S}$. We conclude that
%\begin{equation}\label{eq:tij_reduces_to_zero} 0 = D^{(q+1)(2n-r)-1}\mathcal{S} =  \frac{((q+1)(2n-r)-1)!\,\binom{n}{r}\, (2n-r)}{(q!)^{n-r}\, %\bigl((q+1)!\bigr)^{r} \, \bigl((q+2)!\bigr)^{n - r} }\ t_{i,j}  + E, \end{equation}
%where $E$ belongs to the ideal $(F_0,\ldots , F_{(q+1)(2n-r)})$. 

%Note that $D^{(q+1)(2n-r)-1}\mathcal{S}$ has weight $(q+1)(2n-r)$ and is homogeneous of degree $2n-r$ with respect to the standard grading. The %maximal monomials of this weight and degree are
%\[ y_{q+1}^{2n-r} > y_{q}y_{q+1}^{2n-r-2}y_{q+2} > y_q^2y_{q+1}^{2n-r-4}y_{q+2}^2 > y_{q-1}y_{q+1}^{2n-r-3}y_{q+2}^2\]
%(note that the third monomial does not appear if $n=2$). We show that $\lm(t_{i,j}) = y_{q-1}y_{q+1}^{2n-r-3}y_{q+2}^2$ and that the three bigger %monomials do not occur in (\ref{eq:tij_reduces_to_zero}). This shows that $t_{i,j}$ reduces to zero modulo $\{F_i; i \geq 0\}$. After putting $y_0=0$ %we find that $T_{i,j}$ reduces to zero for $q\geq 2$, as had to be proved. We will deal with the case $q=1$ later. We note that the coefficient of %$F_{(q+1)n}$ in (\ref{eq:tij_reduces_to_zero}) is zero, and hence $y_{q+1}^{2n-r}$ does not occur. To see this, we rewrite $\mathcal{S}$ as 
%\[ \frac{n}{n-r} D(y_0^{n-r}) F_0  - y_0^{n-r} F_1 = 0. \]
%After deriving $(q+1)(2n-r)-1$ times we get 
%\[ \left[\frac{n}{n-r} \binom{(q+1)(2n-r) - 1 }{(q+1)n} - \binom{(q+1)(2n-r) - 1 }{(q+1)n-1}  \right] D^{(q+1)(n-r)}(y_0^{n-r}) \] 
%as coefficient of $F_{(q+1)n}$. A short computation shows that this is zero. The monomial $M:=y_{q}y_{q+1}^{2n-r-2}y_{q+2}$ can then only occur in %$y_{q+2}y_{q+1}^{n-r-1}F_{qn+n-1}$ or $y_{q}y_{q+1}^{n-r-1}F_{qn+n+1}$ but we have just shown that these terms give rise to 
%\[ y_{q+1}^{n-r-1} S(F_{qn+n-1},F_{qn+n+1}) \]
%where $M$ actually gets killed. The next monomial $N:=y_q^2y_{q+1}^{2n-r-4}y_{q+2}^2$ appears in 
%\[ \left\{ \begin{array}{ll} (i) & y_{q+2}y_{q+1}^{n-r-1}F_{qn+n-1} \text{ and }  y_qy_{q+1}^{n-r-1}F_{qn+n+1} , \\
%(ii)  &  y_qy_{q+2}^2y_{q+1}^{n-r-3}F_{qn+n-1}  \text{ and } y_q^2y_{q+2}y_{q+1}^{n-r-3}F_{qn+n+1}, \\
%(iii) & y_{q+2}^2y_{q+1}^{n-r-2}F_{qn+n-2} \text{ and }   y_q^2y_{q+1}^{n-r-2}F_{qn+n+2} . \\ 
% \end{array} \right. \]
%As shown above, the terms in $(iii)$ are part of 
%\[y_{q+1}^{n-r-2}S(F_{qn+n-2},F_{qn+n+2}) \] 
%and $N$ gets killed there. The terms in $(i)$ are part of 
%\begin{equation}\label{eq:some_S_poly} y_{q+1}^{n-r-1}S(F_{qn+n-1},F_{qn+n+1}). \end{equation} 
%Using Proposition~\ref{prop:equations_are_bell_polynomials} and Proposition~\ref{prop:leadingterm_Fi} one may check that the terms containing %$y_q^2y_{q+1}^{n-3}y_{q+2}$ in $F_{qn+n-1}$ and containing $y_qy_{q+1}^{n-3}y_{q+2}^2$ in $F_{qn+n+1}$ are canceled in~(\ref{eq:some_S_poly}). 
%We conclude that there are no terms with $N$ on the left hand side of the equation 
%\[  \frac{((q+1)(2n-r)-1)!\,\binom{n}{r}\, (2n-r)}{(q!)^{n-r}\, \bigl((q+1)!\bigr)^{r} \, \bigl((q+2)!\bigr)^{n - r} }\ t_{i,j}  = -  E, \] 
%and hence $N$ does not appear on the right hand side either. This shows that $N$ gets canceled in the terms in $(ii)$, which are part of $E$, as well. %It remains to study the monomial $K:=y_{q-1}y_{q+1}^{2n-r-3}y_{q+2}^2$. This monomial appears in
%\[ \left\{\begin{array}{l}  y_{q+2}^2y_{q+1}^{n-r-2}F_{(q+1)n-2}, \\ 
% y_{q+2} y_{q+1}^{n-r-1}F_{(q+1)n-1}, \\ 
%y_{q-1}y_{q+2}y_{q+1}^{n-r-2}F_{(q+1)n+1}, \\ 
% y_{q-1}y_{q+1}^{n-r-1}F_{(q+1)n+2}. \end{array} \right. \]
%The first two polynomials occur in $t_{i,j}$ and by the definition of the $S$-polynomial the coefficient of $K$ will be positive there. This is what %we had to show.

%We only have to consider the case $q=1$ now. The biggest monomials of weight $2(2n-r)$ and degree $2n-r$ not containing $y_0$ are
%\[ y_{2}^{2n-r} > y_{1}y_{2}^{2n-r-2}y_{3} > y_1^2y_{2}^{2n-r-4}y_{3}^2 > y_{1}^2y_{2}^{2n-r-3}y_{4}.   \] 
%Now it is not hard to show that this monomial appears in $t_{i,j}$ with negative coefficient. This ends the proof of the lemma.

\begin{Rem}\label{rem:groebner_basis_in_infinitely_many_variables}
We could have avoided to use polynomial rings in countably many variables. In fact the following holds: the leading monomials of
$(f_n,f_{n+1},\ldots,f_{n+m})$ of weight less than or equal to $n+m$ are generated by $\lm(f_i)$, $n\leq i \leq n+m$. In other words: there exists a Gr\"{o}bner basis of $(f_n, \ldots, f_{n+m})$ such that all added elements will be of weight
larger than or equal to $n+m+1$. 
\end{Rem}

\subsection{Computation of the arc Hilbert-Poincar\'{e} series of the $n$-fold point}
Using Gordon's generalization of the Rogers-Ramanujan identity (Theorem~\ref{thm:rogersramanujan}) we immediately obtain an
explicit description of the arc Hilbert-Poincar\'{e} series of the $n$-fold point by a combinatorial interpretation of the leading ideal
$L(I)$ as it was computed
in Lemma~\ref{lem:gbbasis_nfold_point}.

\begin{Thm}\label{thm:hp_nfold_point}
  The Hilbert-Poincar\'{e} series of the focussed arc algebra $J_\infty^0(X)$ of
  the $n$-fold point $X = \{y^n=0\} \subset \mathbb{A}_k^1$ over the origin equals: 
  $$\HP_{J_\infty^0(X)}(t)=\Hh\cdot
  \prod_{\substack{i\geq 1\\ i \equiv 0, n,n+1 \\ \bmod 2n+1}} (1-t^i).$$
  Equivalently, $$\HP_{J_\infty^0(X)}(t)=\prod_{\substack{i\geq 1\\ i\not\equiv 0,n,n+1\\ \bmod 2n+1}} \frac{1}{1-t^i}
  % = \mathbb{H}\cdot \sum_{i=-\infty}^{+\infty}(-1)^it^{\frac{(2n+1)i^2-i}{2}}
  .$$
\end{Thm}

\begin{proof}
  It is a general fact from the theory of Hilbert-Poincar\'{e} series that the Hilbert-Poincar\'{e} series
  of a homogeneous ideal is precisely the Hilbert-Poincar\'{e} series of the leading
  ideal (see Theorem~\ref{thm:hilbert_leading}), i.e.,
  $$\HP_{J_\infty^0(X)}(t)= \HP_{k[y_i;i\geq 1]/L(I)}(t),$$
  where $I$ is as in Lemma~\ref{lem:gbbasis_nfold_point}. By that lemma and Proposition~\ref{prop:leadingterm_Fi} the leading ideal $L(I)$ is generated by monomials of the form $y_{q}^{n-r}y_{q+1}^r$ for $q\geq 1$ and $0\leq r \leq n-1$. Recall
  that the weight of a monomial $y^\alpha=y_{i_1}^{\alpha_1}\cdots
  y_{i_e}^{\alpha_e}$ is precisely $\alpha_1\cdot i_1+ \cdots +
  \alpha_e\cdot i_e$. Thus factoring out $L(I)$ and computing the
  Hilbert-Poincar\'{e} series of the corresponding graded algebra is equivalent to
  counting partitions $(\lambda_1,\ldots,\lambda_s)$ of natural numbers such that $ \lambda_j - \lambda_{j+n-1} > 2$ for all $j$. This is precisely what is counted in Theorem~\ref{thm:gordon}. Hence, we obtain:
  $$\HP_{J_\infty^0(X)}(t)=\prod_{\substack{i\geq 1\\ i\not\equiv 0,n,n+1\\ \bmod 2n+1}} \frac{1}{1-t^i}.$$
  The fact that the right hand side of this equation equals the
  generating series of the number of partitions of $n$ into parts
  which are not congruent to $0,n$ or $n+1$ modulo $2n+1$ is standard in the theory
  of generating series. 
  %The last equality in Theorem~\ref{thm:hp_nfold_point} can be found in standard texts on
  %combinatorics, e.g., Theorem A in Chapter 12 of \cite{rademacher}.
\end{proof}

\subsection{An alternative approach to Rogers-Ramanujan}
In the previous section we used a combinatorial interpretation of the
leading ideal of $I=(f_n,f_{n+1},\ldots)$ to compute the Hilbert-Poincar\'e series of the
corresponding graded algebra. There are commutative algebra methods to
do this as well which yield an alternative approach to the (first) Rogers-Ramanujan
identity. Of course, we consider the case where $n=2$ here, i.e.\ the case of the double point. We will obtain a recursion formula for the
generating functions appearing in the Rogers-Ramanujan identity which has already been considered
by Andrews and Baxter in \cite{andrews_baxter}, though the present
approach gives a natural way to obtain it.\\

Consider the graded algebra $S=k[y_i;i\geq 1]/L(I)$. It is immediate
(see the proof of Theorem~\ref{thm:hp_nfold_point}) that its
Hilbert-Poincar\'e series equals the generating series of the number of partitions of
an integer $n$ without repeated or consecutive parts. Differently, we
compute the Hilbert-Poincar\'e series of $S$ by recursively defining a sequence of
formal power series (generating functions) in $t$ which converges in
the $(t)$-adic topology to the desired Hilbert-Poincar\'e series. We will simply write $k[\geq d]$ for the polynomial ring $k[y_i; i \geq d]$. It will be endowed with the grading $\wt\, y_i = i$. The ideal generated by $y^2_i, y_iy_{i+1}$ for $i\geq d$ in $k[\geq  d]$ will be denoted by $I_d$. We will still write $I_d$ for the ``same'' ideal in $k[\geq d']$ if $d'\leq d$. As usual, if $E$ is an ideal in a ring
$R$ and $f\in R$ then we denote the ideal quotient, i.e., 
\[ \{a\in R\, ;\,  a\cdot f\in E \} \]
by $(E:f)$. \\

Corollary~\ref{cor:hilbert_hilbert_trick_2} implies that $$\HP_{k[\geq
  d]/I_d}(t) = \HP_{k[\geq d]/(I_d,y_d)}(t) + t^d\cdot \HP_{k[\geq
  d]/(I_d:y_d)}(t).$$

Moreover, a quick computation shows the following.

\begin{Prop}
With the notation introduced above we have:
\begin{eqnarray*}
(I_d,y_d)  & = & (y_d,I_{d+1})\\
(I_d:y_d) & = & (y_d,y_{d+1},I_{d+2}).
\end{eqnarray*}
\end{Prop}

This immediately implies 
 $$\HP_{k[\geq  d]/I_d}(t) = \HP_{k[\geq d+1]/I_{d+1}}(t) + t^d\cdot \HP_{k[\geq
  d+2]/I_{d+2}}(t).$$

For simplicity of notation let $h(d)$ stand for $\HP_{k[\geq
  d]/I_d}(t)$. Then:
\begin{equation}\label{eq:recursion}
h(d) = h(d+1) + t^d\cdot h(d+2)
\end{equation}
and:

\begin{samepage}
\begin{Prop}\label{prop:recursion_double}
  For the Hilbert-Poincar\'e series $\HP_{J_\infty^0(X)}(t)=h(1)$ we obtain
$$h(1) = A_d\cdot h(d) + B_{d+1}\cdot h(d+1)$$
for $d\geq 1$ with $A_i,B_i\in k[[t]]$ fulfilling the following
recursion
\begin{eqnarray*}
  A_d & = & A_{d-1} + B_d\\
  B_{d+1} & = & A_{d-1}\cdot t^{d-1}
\end{eqnarray*}
with initial conditions $A_1=A_2=1$ and $B_2=0, B_3=t$.
\end{Prop}
\end{samepage}

\begin{proof}
  By the discussion above $h(1)$ equals $ h(2) + t\cdot h(3)$; hence,
  $A_1=A_2=1$ and $B_2=0$, $B_3=t$. Assume now that $$h(1) = A_d\cdot
  h(d) + B_{d+1}\cdot h(d+1)$$ holds for some $d\geq 2$. By
  equation (\ref{eq:recursion}) substituting for $h(d)$ yields
  \begin{eqnarray*}
    h(1) & = & A_d\cdot (h(d+1) + t^d\cdot h(d+2)) + B_{d+1}\cdot h(d+1)\\
    & = & (A_d + B_{d+1})\cdot h(d+1) + (A_d\cdot t^d)\cdot h(d+2)
  \end{eqnarray*}
  from which the assertion follows.
\end{proof}

If $(s_d)_{d\in \N}$ is a sequence of formal power series $s_d\in
k[[t]]$ we will denote by $\lim s_d$ its limit -- if it exists -- in
the $(t)$-adic topology. Since $\ord\, B_d \geq d-2$ it is immediate
that both $\lim A_d$ and $\lim B_d$ exist, in fact: $\lim B_d=0$ and
$$h(1) = \lim A_d.$$
The recursion from Proposition \ref{prop:recursion_double} can easily
be simplified. We obtain:

\begin{Cor}\label{cor:recursion_rogersramanujan}
  With the above introduced notation
  $\HP_{J_\infty^0(X)}(t)=\lim A_d$ where $A_d$ fulfills $$A_d = A_{d-1} + t^{d-2}\cdot A_{d-2}$$
  with initial conditions $A_1=A_2=1$.
\end{Cor}

The recursion appearing in this corollary is well-known since Andrews and Baxter \cite{andrews_baxter}. Its limit is
precisely the infinite product
\[ \prod_{\substack{i\geq 1\\ i\equiv 1, 4 \bmod 5}} \frac{1}{1-t^i}, \]
i.e., the generating series of the number of partitions with parts equal to $1$ or $4$ modulo $5$. Note, that our
construction gives the generating series $G_i$ defined in the paper by Andrews and Baxter an interpretation as
Hilbert-Poincar\'e series of the quotients $k[\geq i]/I_i$. This immediately implies that the series $G_i$ are of the
form $G_i = 1 + \sum_{j\geq i} G_{ij} t^j$ (this observation was called an `empirical hypothesis' by Andrews and Baxter).

%If we set $S_d=A_{d+2}$ then
%$h(1)$ is also the limit of $(S_d)$ which equals the sum $$S_d=
%\sum_{i=0}^n T_i, T_i=S_{i-2}\cdot t^i$$ where $T_0=1$,$T_1=t$ which
%is according to OEIS a recursion fulfilled by the Rogers-Ramanujan
%function. 


